\section{Доказательство оптимальности}
% Доказательство оптимальности симплекс-метода
%
% Описанную задачу можно решить симплекс-методом, т.к.:
%
% \begin{enumerate}
%     \item Целевая функция и все ограничения задачи \textbf{линейны}, то есть заданы в виде линейных комбинаций переменных;
%     \item Все переменные задачи \textbf{неотрицательны};
%     \item Область допустимых решений \textbf{ограничена}. Неотрицательность переменных обеспечивает ограничение снизу, а дополнительные условия из таблицы определяют верхнюю границу;
%     \item Задача \textbf{разрешима}. Существует хотя бы одно допустимое решение, удовлетворяющее всем ограничениям: $x_1 = 1,\quad x_2 = 0,\quad x_3 = 0,\quad \overline{x_4} = 0,\quad \overline{x_5} = 0,\quad \overline{\overline{x_4}} = 0,\quad \overline{\overline{x_5}} = 0,\quad s_1 = 5,\quad s_2 = 0$;
%     \item Задача имеет \textbf{конечный оптимум}: поскольку область допустимых решений ограничена, а целевая функция линейна, оптимальное значение достигается в одной из вершин многогранника допустимых решений.
% \end{enumerate}
%
% Таким образом, симплекс-метод гарантированно найдет оптимальное решение, перебирая вершины многогранника допустимых решений. Кроме того, по теореме двойственности, \textbf{задача, двойственная к данной, также имеет решение}, а значения функций цели в оптимумах совпадут.

Для доказательства того, что симплекс-метод находит оптимальное решение данной задачи линейного программирования, необходимо показать выполнение условий теоремы оптимальности симплекс-метода.

\subsection*{1. Формализация задачи}
Задача линейного программирования записана в каноническом виде:

$$Z = c^T x \rightarrow \min$$

при ограничениях:

$$Ax = b, \quad x \geq 0.$$

где:
\begin{itemize}
    \item $A \in \mathbb{R}^{m \times n}$ – матрица коэффициентов ограничений,
    \item $b \in \mathbb{R}^m$ – вектор правых частей,
    \item $c \in \mathbb{R}^n$ – вектор коэффициентов целевой функции.
\end{itemize}

\subsection*{2. Достаточные условия оптимальности}
Симплекс-метод гарантирует нахождение оптимального решения, если выполняются следующие условия:

\begin{enumerate}
    \item Существование базисного допустимого решения
    \begin{itemize}
        \item В начальном базисе присутствуют $m$ линейно независимых столбцов матрицы $A$, соответствующих базисным переменным.
        \item В данной задаче присутствуют уравнения (равенства), поэтому базисное решение можно найти. Если его нет сразу, оно строится с помощью искусственных переменных.
    \end{itemize}

    \item Ограниченность множества допустимых решений
    \begin{itemize}
        \item Если оптимальное решение существует, симплекс-метод его найдёт, так как в данной задаче переменные $x_i \geq 0$, что исключает случай неограниченной допустимой области.
    \end{itemize}

%     \item Невырожденность базисного решения или обработка вырожденности
%     \begin{itemize}
%         \item В случае, если базисное решение вырождено (некоторые базисные переменные равны нулю), возможно зацикливание. Однако в симплекс-методе применяются правила (например, правило Бленда), позволяющие избежать этой проблемы.
%     \end{itemize}

    \item Остановка на оптимальном базисе
    \begin{itemize}
        \item Симплекс-метод завершается, когда векторы улучшения (оценки в таблице симплекс-метода) удовлетворяют условию:
        $$c_N - c_B B^{-1} A_N \leq 0.$$
        Это означает, что дальнейшее улучшение целевой функции невозможно.
    \end{itemize}
\end{enumerate}

\subsection*{3. Теорема двойственности и оптимальность решения}
Дополнительно можно доказать оптимальность с помощью двойственной задачи:

$$y^T A \geq c^T, \quad y^T b = Z^*.$$

где $y$ – двойственные переменные.

\begin{itemize}
    \item Если найдены такие $y$, что выполняется это равенство, значит, достигнута оптимальность.
    \item В данной задаче двойственная задача строится стандартным способом, и её решение совпадает с прямой задачей, что подтверждает нахождение глобального максимума.
\end{itemize}

\subsection*{Вывод}
Поскольку в данной задаче:
\begin{enumerate}
    \item Допустимая область ограничена и не пустая (из-за наличия неотрицательных переменных и уравнений),
    \item Целевая функция и ограничения линейны,
    \item Симплекс-метод корректно завершается, когда оценки $\Delta_j \leq 0$ для всех небазисных переменных,
    \item Двойственная теорема подтверждает совпадение значений оптимума,
\end{enumerate}
